% Loadshift -- Ignition Hacks V.7 demo video
% Three parts: A) what to record, B) what to say, C) the voiceover script.
% Picture and audio are recorded separately; each fits inside 3:00 on its own.
% Compiles with plain pdflatex.

\documentclass[10pt]{article}

\usepackage[T1]{fontenc}
\usepackage[utf8]{inputenc}
\usepackage[letterpaper,margin=0.85in]{geometry}
\usepackage{xcolor}
\usepackage{booktabs}
\usepackage{longtable}
\usepackage{array}
\usepackage{enumitem}
\usepackage{titlesec}
\usepackage{parskip}

\definecolor{spruce}{HTML}{1F6F53}
\definecolor{amber}{HTML}{B4761A}
\definecolor{ink3}{HTML}{6B7280}

\titleformat{\section}{\large\bfseries\color{spruce}}{\thesection.}{0.5em}{}
\titleformat{\subsection}{\normalsize\bfseries}{}{0em}{}
\setlist[itemize]{leftmargin=1.2em,itemsep=1pt,topsep=2pt}
\setlist[enumerate]{leftmargin=1.4em,itemsep=1pt,topsep=2pt}

\newcommand{\cotwo}{CO\textsubscript{2}}
\newcommand{\dur}[1]{{\small\ttfamily #1}}
\newcommand{\prize}[1]{{\small\color{spruce}\textbf{[#1]}}}
\newcommand{\cut}{{\small\color{amber}\textbf{[cuttable]}}}
\newcommand{\tc}[1]{{\ttfamily\bfseries #1}}

\newcolumntype{R}[1]{>{\raggedright\arraybackslash}p{#1}}

\begin{document}

\begin{center}
{\LARGE\bfseries Loadshift --- demo video}\\[3pt]
{\small Ignition Hacks V.7, Environmental track \quad$\cdot$\quad hard cap \textbf{3:00}}\\[2pt]
{\small\color{ink3} Targeting: Best Solo Hack \quad$\cdot$\quad Best Use of Render \quad$\cdot$\quad Best Overall}
\end{center}

\vspace{2pt}
\hrule
\vspace{6pt}

\noindent
\textbf{How to use this.} Part~A is the capture pass --- record it silently, in order,
and it lands at \textbf{2:52} of footage. Part~B is every claim worth making, ranked, so
you know what to drop if you run long. Part~C is the voiceover you read over the
finished cut: \textbf{411 words}, which is \textbf{2:50} of speech at 145 words per
minute. Picture and audio do not need to line up beat for beat --- they are the same
length by construction, so you can nudge either one in the edit.

\medskip
\noindent
\textbf{Weighting.} Because Best Use of Render is a target, the Render architecture is
not a closing footnote. It is the single largest block of the voiceover --- a third of
all the words --- and it gets seven of the twenty-one shots. The ML method is present
but compressed.

% =====================================================================
\section{Part A --- The shot list (what to record)}

Record against the \textbf{live site}, \texttt{loadshift-web.onrender.com}, so the URL
bar is itself evidence the thing is deployed.

\subsection{Before you hit record}

\begin{itemize}
  \item \textbf{Warm the instance.} Load the site a few minutes early. Render's starter
        tier cold-starts, and the first paint after a sleep is slow.
  \item \textbf{Pre-run the sample analysis once.} The generated report is cached in
        Key Value for 24h, keyed by a hash of the statistics, so running it before the
        take makes it return fast during the take.
  \item \textbf{Clean browser profile.} No bookmarks bar, no extensions, no other tabs.
        Zoom at 100\%.
  \item \textbf{The page refetches every 5 minutes.} A number can change mid-take. That
        is not a bug, but keep takes short so it does not change mid-shot.
  \item Have \texttt{api/samples/demo\_upload\_greenbutton.xml} somewhere obvious for
        the file dialog.
\end{itemize}

\subsection{Site footage --- one continuous scroll \hfill \normalfont 1:54}

\begin{longtable}{@{}R{0.4cm} R{3.2cm} R{1.0cm} R{9.4cm}@{}}
\toprule
\textbf{\#} & \textbf{Shot} & \textbf{Dur} & \textbf{What to do on screen} \\
\midrule
\endhead
1 & Hero + the two numbers & \dur{12s} &
Load the page. Hold on the stacked pair: \emph{What your load causes} beside \emph{What
the news reports}, and the line about an added load being 1.6$\times$ dirtier than the
average suggests. This frame is the whole thesis --- give it room. \\
\addlinespace
2 & Live badge & \dur{5s} &
Pause on the pulsing green pill: \texttt{Live $\cdot$ Ontario grid $\cdot$ <time>}. It
proves the data is current, not a screenshot. \\
\addlinespace
3 & The Day Band & \dur{10s} &
The 24-hour strip, colour-scaled by forecast marginal intensity, with the
\texttt{cleanest} flag above it. Hover two or three cells so the g/kWh tooltip pops.
Show the dimmed night hours. \\
\addlinespace
4 & Live fuel mix & \dur{4s} &
The monospace row underneath: nuclear / hydro / gas / wind, in MW, right now. \\
\addlinespace
5 & Pick the load & \dur{8s} &
Scroll to \emph{Schedule}. Click Dryer, then EV charge --- the answers recompute. Type a
nameplate wattage into the optional box. \\
\addlinespace
6 & Awake-hours slider & \dur{10s} &
Drag both thumbs of the wake/bed picker. The day band re-shades and the answers update
live. Drag it past midnight once so the awake span wraps into two segments. \\
\addlinespace
7 & The two answers & \dur{12s} &
Hold on both cards: \emph{While you're up} and \emph{Best overall $\cdot$ while you
sleep}. Percent saved, grams, the ``this hour vs the worst hour'' dollar line, the
km-of-driving equivalent, and the Ultra-Low Overnight note. \\
\addlinespace
8 & Markers on the band & \dur{5s} &
The large day band below, now flagged with \texttt{while you're up} and
\texttt{overall}. \\
\addlinespace
9 & Green Button upload & \dur{8s} &
Scroll to \emph{Your data}. Click \textbf{Upload Green Button XML} and pick the sample
from the file dialog. \textbf{Record the dialog} --- it shows this takes a real file,
not a demo button. \\
\addlinespace
10 & The six stat tiles & \dur{10s} &
Emissions last year, recoverable kg and \%, bill savings, timing score, evening-peak
share, and the rate-plan verdict. \\
\addlinespace
11 & Usage charts & \dur{7s} &
Both bar charts. Hover one hour and one month so the tooltips appear. \\
\addlinespace
12 & Generated report & \dur{10s} &
The \emph{Your energy report} card with the \texttt{AI-generated $\cdot$ Gemini} chip:
summary plus numbered recommendations. Then type one follow-up --- ``Why is my evening
usage a problem?'' --- and show the answer. \\
\addlinespace
13 & Model metrics & \dur{8s} &
\emph{Method} section: forecast error, naive baseline, improvement \%, test $R^2$. These
are read live from the model card, not typed into the page. \\
\addlinespace
14 & Known limitation & \dur{5s} &
The amber card, computed live from the model artifact: the model covers about 78\% of
Ontario's extra demand. Putting the weakness on camera counts in your favour, not
against you. \\
\bottomrule
\end{longtable}

\subsection{Architecture and code \hfill \normalfont 0:58 --- screen capture or stills}

\begin{longtable}{@{}R{0.4cm} R{3.2cm} R{1.0cm} R{9.4cm}@{}}
\toprule
\textbf{\#} & \textbf{Shot} & \textbf{Dur} & \textbf{What to do on screen} \\
\midrule
\endhead
15 & Render dashboard & \dur{12s} &
\textbf{The most valuable frame in the whole video.} The Blueprint view with all four
resources visible together: \texttt{loadshift-web}, \texttt{loadshift-api},
\texttt{loadshift-refresh}, \texttt{loadshift-cache}. Make sure all four are in one
screenshot. \\
\addlinespace
16 & \texttt{render.yaml} & \dur{10s} &
Slow scroll past the four \texttt{type:} blocks --- \texttt{keyvalue}, \texttt{web},
\texttt{cron}, \texttt{web}. Land on the \texttt{fromService} wiring that injects
\texttt{KV\_URL}, and on \texttt{maxmemoryPolicy: noeviction}. \\
\addlinespace
17 & Cron run history & \dur{8s} &
\texttt{loadshift-refresh} job history: a column of hourly successes at seven past. If
the log is legible, the line where it publishes the forecast. \\
\addlinespace
18 & \texttt{/api/health} & \dur{6s} &
Raw JSON in a browser tab. The three fields that matter: \texttt{kv\_ok: true},
\texttt{cache\_backend: render-key-value}, \texttt{refresher: cron}. \\
\addlinespace
19 & Architecture diagram & \dur{8s} &
The topology block from \texttt{DEPLOY.md}: browser $\rightarrow$ web $\rightarrow$ api
$\rightarrow$ key value $\leftarrow$ cron. Easiest as a still; redraw it if you want it
prettier. \\
\addlinespace
20 & The ML core \cut & \dur{8s} &
Still or slow scroll: the marginal-emissions regression in
\texttt{api/loadshift/mef.py}, or the 14-feature list in
\texttt{api/loadshift/model.py}. \\
\addlinespace
21 & Tests passing \cut & \dur{6s} &
\texttt{pytest tests/ -q} finishing at 63 passed. \\
\bottomrule
\end{longtable}

\subsection{Timing}

\begin{tabular}{@{}lr@{}}
\toprule
Site footage (shots 1--14) & 1:54 \\
Architecture and code (shots 15--21) & 0:58 \\
\midrule
\textbf{Total} & \textbf{2:52} \\
Without shots 20--21 & 2:38 \\
\bottomrule
\end{tabular}

\medskip
\noindent\textbf{Drop in this order if you run long:} 21, 20, 4, 8, 11. Never drop 15
or 16 --- those two are the Render prize. Never drop 1 or 7 --- those two are the
product. Every duration above is a ceiling, not a target; coming in under is free
slack for the merge.

% =====================================================================
\newpage
\section{Part B --- Talking points (everything worth saying)}

Ranked. Each group is tagged with the prize it serves. Part~C uses the top of this list;
the rest is here for the Devpost writeup and for judge questions afterwards.

\subsection{1. The thesis \prize{Overall}}
\begin{itemize}
  \item Ontario's grid averages around 100 g\cotwo{}/kWh and looks clean --- mostly
        nuclear and hydro. So ``does it matter when I run the dryer?'' sounds like a
        non-question here.
  \item It matters, because the average is the wrong number. Nuclear runs flat no
        matter what you do. The plant that ramps up to serve one more kWh is almost
        always natural gas, at roughly 490 g.
  \item Across the labelled history: mean marginal 227 g against a mean average of
        103 g. A 2.2$\times$ gap --- the entire thesis in one number.
  \item And unlike the average, the marginal number swings about 50\% within a single
        day. That swing is the lever.
  \item Nobody publishes a marginal forecast for Ontario households, so a lever
        available to every home on the grid goes unused.
\end{itemize}

\subsection{2. What it does \prize{Overall}}
\begin{itemize}
  \item Forecasts marginal intensity 24 hours ahead and names the cleanest hour to run
        a deferrable load --- dryer, dishwasher, washer, EV charger.
  \item \textbf{Two} answers, not one: the best window while you're awake, and the best
        window overall. You have to be up to start a load, but it can finish while you
        sleep --- that is what a delay-start timer is for.
  \item Percent saved is the headline because it is \emph{appliance-independent}: the
        kWh cancels out, so the number is identical for an efficient dryer and an
        ancient one. It depends only on the grid. Chasing a precise gram figure for an
        appliance you cannot see would be fake precision.
  \item Dollars are costed against the OEB Regulated Price Plan --- both Time-of-Use
        and Ultra-Low Overnight --- so the tool also tells you whether you are on the
        wrong rate plan.
\end{itemize}

\subsection{3. Real meter data \prize{Overall}}
\begin{itemize}
  \item Green Button, under O.~Reg.~633/21: every Ontario utility must hand you your
        own smart-meter interval history. Upload the XML and get a year of analysis.
  \item The bundled EV household: 9,563 kWh a year, 2,421.81 kg \cotwo{} actual against
        2,228.02 kg optimally shifted. \textbf{193.79 kg and \$82.93 recoverable purely
        by moving when the load runs.}
  \item Also returned: a timing score, the evening-peak share, and a verdict on the
        rate plan. Files are analysed in memory and never stored.
  \item A Gemini-generated report reads the statistics back in plain language and
        answers follow-up questions --- answered \emph{only} from those statistics.
\end{itemize}

\subsection{4. Method, stated honestly \prize{Overall + Solo}}
\begin{itemize}
  \item No public marginal-emissions feed exists for Ontario, so I derived the label:
        hour-over-hour deltas of each fuel's output and of demand, keeping only hours
        where demand rose, bucketed by season and hour-block, regressed within
        overlapping net-demand windows. Following Siler-Evans, Azevedo and Morgan
        (2012), \emph{Environ.\ Sci.\ Technol.}\ 46(9).
  \item Marginal intensity is the sum of each fuel's slope times its IPCC AR5 lifecycle
        emission factor.
  \item Forecast: LightGBM, 24h ahead, 14 features, every one causal at that horizon
        --- calendar, forecast weather, and lags of 24h or more. No leakage. 23,112
        hourly rows, 2.6 years, from 2024-01-01.
  \item \textbf{I wrote the seasonal-naive baseline first, deliberately, so the result
        would mean something.} On a chronological 60-day holdout: MAE 13.51 g against
        23.21 g. 41.8\% better, $R^2$ 0.555. ``13.5 g'' alone means nothing; ``13.5
        against 23.2'' is a result.
  \item Externally validated: my independently computed average intensity tracks
        Electricity Maps' published Ontario figure to within 2.1 g/kWh, bias +1.8 g,
        correlation 0.91. Validated against them once; never built on them.
  \item Stated limitation --- first line of the limitations doc and rendered live in the
        UI: the regression explains about 78\% of a marginal Ontario kWh. The rest is
        intertie imports and exports, which I do not observe.
\end{itemize}

\subsection{5. Render --- the headline block \prize{Best Use of Render}}
\begin{itemize}
  \item \textbf{Four resources, one Blueprint, one \texttt{git push}.}
        \texttt{loadshift-web} (Next.js), \texttt{loadshift-api} (FastAPI),
        \texttt{loadshift-refresh} (hourly cron), \texttt{loadshift-cache} (Key Value,
        Valkey 8). All declared in \texttt{render.yaml} at the repo root --- infra as
        code, nothing clicked into the dashboard. \texttt{KV\_URL} is wired by
        \texttt{fromService}, so no connection string is ever copied by hand.
  \item \textbf{Private network.} The browser only ever calls the web service's own
        origin; \texttt{/api/*} is proxied to the API's \emph{private} hostname. No CORS
        preflight, no \texttt{onrender.com} API URL baked into the frontend bundle, and
        API traffic never leaves Render's network. It is a route handler rather than a
        \texttt{next.config} rewrite on purpose --- rewrites freeze at build time, which
        would mean the API's hostname had to exist before the frontend could build.
  \item \textbf{The ephemeral-filesystem lesson.} Render containers destroy the disk on
        every deploy, and I learned that the expensive way. Four things lived in the web
        container: the forecast cache, the last-known-good weather, the shared Gemini
        key budget, and the generated-report cache. Ordinary deploys wiped all four ---
        and the daily API cap was multiplying by the number of running instances. All
        four now live in Render Key Value, the only tier that survives a deploy. Set to
        \texttt{noeviction}, because the forecast must never be dropped to make room for
        a cache entry.
  \item \textbf{The topology enforces the rule.} The cron job is the only process that
        ever loads LightGBM. ``Never run the model on a request path'' is therefore
        enforced by architecture, not by discipline --- the web service \emph{cannot}
        run the model, because it never imports it. Redeploy the API and the very first
        request returns a warm forecast.
  \item \textbf{Liveness is not readiness.} \texttt{/api/health} is the
        \texttt{healthCheckPath} and always returns 200: an instance still waiting on
        its first forecast is alive, and failing its own health check would make Render
        roll back a perfectly good deploy. \texttt{/api/ready} is the readiness probe
        and 503s when there is genuinely nothing to serve.
  \item \textbf{The supervisor.} If no refresh job has reported in for two hours, the
        web service rebuilds the forecast itself and reports
        \texttt{refresher: web-fallback}. It is a safety net, not a second scheduler ---
        when the cron is healthy it does nothing. It exists because I once shipped the
        cron-based code \emph{before} the cron job existed, and the forecast froze and
        aged silently with nothing scheduled to replace it.
  \item \textbf{Never an error page.} IESO or the weather API down: last good forecast,
        flagged stale, with the time it was built. A cron run fails: same. Key Value
        unreachable: every helper returns \texttt{None}, each service falls back to its
        in-process path, and health reports
        \texttt{cache\_backend: in-process (kv unreachable)}. Every failure mode
        degrades to something a visitor can still use.
\end{itemize}

\subsection{6. Solo scope \prize{Best Solo Hack}}
\begin{itemize}
  \item One person, about 22 hours, roughly 5,400 lines, real smart-meter data working
        end to end.
  \item 63 tests across 6 files --- the IESO parsers and hour conventions, TOU and ULO
        pricing, the weather fallback chain, the insight cache, and a file that
        deliberately simulates the durable cache being unreachable. One of those test
        files exists purely because of the outage described above.
  \item A public 27-item document of everything the project assumes and everything it
        cannot do. Naming the weaknesses is part of the product.
  \item No API keys on the main path. Every data source is public and keyless.
\end{itemize}

\subsection{7. Sharp edges I hit \prize{Solo} \cut}
\begin{itemize}
  \item Two IESO XML reports use two different namespaces. Hours are numbered 1--24,
        not 0--23. Reports use fixed EST year-round, which I only trusted after checking
        that the daylight-saving transition days have exactly 24 rows. One hour is
        missing from the record entirely, at the Market Renewal go-live.
  \item A single cp1252-encoded byte in a Python comment broke the Render build.
  \item \texttt{Expect: 100-continue} was being rejected on Green Button uploads.
  \item The insight cache kept missing because Python serialises \texttt{7036.0} where
        JavaScript sends \texttt{7036}, so identical statistics hashed differently.
\end{itemize}

\subsection{8. What's next \cut}
\begin{itemize}
  \item Model intertie flows and close the 22\% gap; per-plant heat rates instead of
        IPCC medians for the gas fleet.
  \item Utility OAuth, so Green Button pulls automatically instead of being uploaded.
  \item Smart plugs, EV chargers and thermostats, so Loadshift \emph{executes} the shift
        rather than recommending it.
  \item Other grid operators --- the method is not Ontario-specific, only the data
        plumbing is.
\end{itemize}

% =====================================================================
\newpage
\section{Part C --- The voiceover script}

\noindent
\textbf{411 words.} At 145 words per minute that is \textbf{2:50} of speech; leave the
remaining ten seconds for the breaths marked below. Read it unhurried --- if you rush
to fit, it sounds rushed. Two fragments are marked \cut{} and buy back about eight
seconds if a take runs long.

\medskip
\noindent
Numbers are written the way you say them, not the way you would type them. The
\emph{On screen} line is a hint for the edit, not a hard sync point.

\vspace{4pt}
\hrule
\vspace{6pt}

%<vo>
\subsection{\tc{0:00 -- 0:17} \quad Hook \hfill {\small\color{ink3} 42 words}}
\emph{On screen: shots 1--2, the hero and the two numbers.}

\medskip
Ontario's grid looks clean. About a hundred grams of \cotwo{} per kilowatt-hour, mostly
nuclear and hydro. But the average is the wrong number. Nuclear runs flat. The plant
that ramps up for your extra load is gas, at four hundred and ninety.

\vspace{8pt}
\subsection{\tc{0:17 -- 0:38} \quad Now \hfill {\small\color{ink3} 49 words}}
\emph{On screen: shots 2--4, the live numbers, the day band, the fuel mix.}

\medskip
Loadshift forecasts that marginal number twenty-four hours ahead. This is right now:
what your load actually causes, beside what the news reports. An added kilowatt-hour is
one point six times dirtier than the average suggests. And here's the next twenty-four
hours as one band --- green is clean, amber isn't.

\vspace{8pt}
\subsection{\tc{0:38 -- 0:59} \quad Schedule \hfill {\small\color{ink3} 52 words}}
\emph{On screen: shots 5--8, appliance, slider, the two answer cards.}

\medskip
Pick a load, set the hours you're awake, and you get two answers: the best window while
you're up, and the best one overall. Percent saved is the headline, because the
kilowatt-hours cancel out. It's the same number for a new dryer and a twenty-year-old
one. The dollars use Ontario's regulated rates.

\vspace{8pt}
\subsection{\tc{0:59 -- 1:19} \quad Your data \hfill {\small\color{ink3} 47 words}}
\emph{On screen: shots 9--12, the upload dialog, the tiles, the report.}

\medskip
This is Green Button data --- the history every Ontario utility must give you. Upload it
and you get a year back. This household emitted two thousand four hundred kilograms of
\cotwo{}. Moving the same loads to cleaner hours saves a hundred and ninety-four kilos,
and eighty-three dollars.

\vspace{8pt}
\subsection{\tc{1:19 -- 1:40} \quad Method \hfill {\small\color{ink3} 51 words}}
\emph{On screen: shots 13--14 and 20, the metric tiles and the limitation card.}

\medskip
Nobody publishes a marginal forecast for Ontario, so I derived one, from hour-over-hour
fuel deltas on the hours demand rose. Then LightGBM, on twenty-three thousand hours of
history. I wrote the naive baseline first, deliberately, and the model beats it by
forty-two percent. \cut{} The page also states what it can't see.

\vspace{8pt}
\subsection{\tc{1:40 -- 2:38} \quad Render \hfill {\small\color{ink3} 140 words --- the big one}}
\emph{On screen: shots 15--19 --- dashboard, \texttt{render.yaml}, cron history, health
JSON, diagram.}

\medskip
Loadshift is four Render resources in one Blueprint: a Next.js front end, a FastAPI
service, a Key Value store, and an hourly cron job. One git push deploys all of it.

\medskip
The browser only ever talks to the web service. Slash-A-P-I is proxied over Render's
private network, so API traffic never leaves it.

\medskip
Render's filesystem is ephemeral, and I learned that the hard way --- four caches lived
on disk, and every deploy destroyed them. They're in Key Value now, the only tier that
survives a deploy.

\medskip
And the cron job is the only process that ever loads the model. The web service doesn't
even import LightGBM, so it can't run the model on a request path. If an upstream fails,
you get the last good forecast, flagged stale. Nothing here ever shows an error page.

\vspace{8pt}
\subsection{\tc{2:38 -- 2:50} \quad Close \hfill {\small\color{ink3} 30 words}}
\emph{On screen: shot 21 if you kept it, then back to the hero.}

\medskip
Solo build, about twenty-two hours, sixty-three tests, and a public document of
twenty-seven things it assumes and can't do. Ontario already has the clean hours.
Loadshift just tells you when.
%</vo>

\vspace{10pt}
\hrule
\vspace{6pt}

\subsection*{Delivery notes}
\begin{itemize}
  \item The Render block is four short paragraphs on purpose. Take a real breath between
        them --- it is the densest thing you say, and the pauses are what make it land.
  \item Say ``slash-A-P-I'', not ``slash api''.
  \item ``LightGBM'' is read as ``light G B M''.
  \item In the hook, the beat is after \emph{wrong number}. Let it sit for half a second
        before ``Nuclear runs flat''.
  \item The Gemini report is a visual beat only --- shot 12 shows it while the Your data
        block is playing, so nothing in the script has to name it.
  \item If a take runs past 2:55, cut the two \cut{} fragments first.
\end{itemize}

\subsection*{Number check}
Every figure quoted here traces to the repo: the model metrics to
\texttt{api/artifacts/model\_card.json}, the Green Button figures to the bundled sample,
the architecture to \texttt{render.yaml} and \texttt{DEPLOY.md}. Two numbers are live
and can move before you record. The metrics in shot 13 are read from the model card, so
if you retrain, re-check the forty-two percent in the Method block against what the page
displays. The 1.6$\times$ multiplier in the Now block is whatever the grid is doing at
that moment --- read the hero before you commit that take.

\end{document}
